% =====================================================================
%  Chapitre 4 — Interfaces
% =====================================================================
\chapter{Présentation des interfaces}
\label{chap:interfaces}

Ce chapitre parcourt les quatorze écrans de l'application dans l'ordre où une
démonstration les enchaîne naturellement. Chaque écran est décrit puis accompagné d'un
encadré \emph{Présentation orale suggérée}, directement réutilisable en soutenance.

\begin{remarque}[Insertion des captures]
Les captures sont appelées depuis le répertoire \code{rapport/images/}. Tant qu'un
fichier n'y figure pas, le document compile et affiche un cadre nommant l'image
attendue. Il suffit de déposer le fichier au nom indiqué pour que la capture apparaisse.
\end{remarque}

% =====================================================================
\section{Connexion}

L'écran de connexion est construit en deux panneaux. À gauche, sur fond vert sombre, la
promesse du produit et les quatre temps de la démarche, numérotés de
\emph{01~—~Diagnostiquer} à \emph{04~—~Piloter la valeur}. À droite, sur fond clair, le
formulaire : adresse électronique professionnelle, mot de passe avec bascule
d'affichage, et bouton de connexion doté d'un état de chargement.

Quatre états d'erreur distincts sont rendus en clair plutôt qu'en code HTTP :
identifiants invalides, compte désactivé, saisie incomplète et serveur injoignable.

Aucun identifiant de démonstration n'est écrit dans le composant React. Lorsque le
serveur tourne en mode développement, un panneau discret propose la liste des comptes,
obtenue d'un endpoint qui répond 404 en production. En production, ce panneau n'existe
donc tout simplement pas.

\capture{01-connexion.png}{Écran de connexion}{0.92}

\begin{oral}
Voici la page d'entrée de la plateforme. Le panneau de gauche résume la démarche en
quatre temps, du diagnostic au pilotage de la valeur. Le formulaire est volontairement
sobre. Un point de sécurité mérite d'être signalé : aucun compte de démonstration n'est
écrit en dur dans le code de l'interface. La liste que vous voyez ici vient d'un
endpoint qui n'existe qu'en mode développement et qui répond « non trouvé » en
production. Je me connecte avec le compte de la Direction Générale.
\end{oral}

% =====================================================================
\section{Tableau de bord Direction}

C'est l'écran d'accueil des utilisateurs porteurs de la permission
\code{dashboard.direction}. Il est composé de quatre indicateurs de tête — projets
actifs, investissement, ROI du portefeuille et risques critiques — chacun assorti d'une
précision de contexte calculée, par exemple « engagé sur 3,1~M\texteuro{} budgétés ».

Au-dessous, la matrice Impact~$\times$~Faisabilité positionne chaque opportunité, la
taille de la bulle représentant le bénéfice annuel attendu. À côté, la répartition du
portefeuille par phase TOGAF, de la qualification à la production. Enfin, un tableau des
premières priorités, cliquable jusqu'à la fiche projet.

\capture{02-dashboard-direction.png}{Tableau de bord Direction}{0.95}

\begin{oral}
Voici la vue de la Direction Générale. Ces quatre chiffres ne sont pas saisis : ils sont
agrégés à la demande depuis les projets, les coûts et les bénéfices. Le ROI de
portefeuille est le rapport entre les gains cumulés sur trois ans et le coût total de
possession. La matrice croise l'impact métier et la faisabilité ; la taille de chaque
bulle est le bénéfice annuel attendu. Le quadrant en haut à droite est ce que
l'organisation devrait lancer en premier. Toutes les lignes du tableau sont cliquables
et mènent à la fiche projet correspondante.
\end{oral}

% =====================================================================
\section{Tableau de bord Département}

Le même emplacement sert une vue différente à qui porte
\code{dashboard.department} : opportunités du périmètre, projets actifs, gains annuels
attendus et risques ouverts. Le périmètre est celui des départements de rattachement de
l'utilisateur, jamais l'organisation entière.

Suivent les opportunités du périmètre classées par score, la répartition par phase, les
projets en cours avec leur avancement et leur prochaine gate, les risques classés par
criticité et les indicateurs métier avec la part de l'écart déjà captée entre la
référence et la cible.

\capture{03-dashboard-departement.png}{Tableau de bord Département}{0.95}

\begin{oral}
Même URL, écran différent. Je suis maintenant connecté en responsable de département.
La plateforme ne se contente pas de masquer des informations : le backend filtre les
données sur les départements de rattachement de l'utilisateur. Si j'appelais l'API
directement pour lire une opportunité d'un autre département, je recevrais un 404, pas
un 403 — répondre 403 confirmerait que la ressource existe.
\end{oral}

% =====================================================================
\section{Diagnostic AS-IS et readiness}

L'écran affiche la campagne en cours, son implantation et son taux de complétion. Une
jauge donne le readiness global sur 100 et le niveau de maturité associé, de
« Inexistant » à « Optimisé ». À droite, l'échelle de maturité détaillée, puis deux
bandeaux : les dimensions bloquantes, en ambre, et les points de vigilance.

Deux colonnes séparent le readiness technique du readiness métier et organisationnel,
chacune assortie de sa moyenne sur~5. La dernière carte, « Preuves collectées »,
présente chaque dimension avec sa note, son poids, sa confiance et la preuve qui
l'étaye.

Deux actions sont proposées : exporter la campagne au format CSV et poursuivre le
diagnostic. Le second bouton s'adapte au droit de l'utilisateur : il devient « Voir le
questionnaire » pour qui n'a pas \code{diagnostic.answer}.

\capture{04-diagnostic.png}{Diagnostic AS-IS et calcul du readiness}{0.95}

\begin{oral}
Le readiness est le point de départ de la démarche. Chaque dimension est notée de 0 à 5,
pondérée par son importance et par la confiance que le répondant accorde à sa propre
réponse. Cette confiance intervient au numérateur et au dénominateur : une réponse mal
étayée ne fausse le score ni vers le haut ni vers le bas. Les deux dimensions en ambre
sont bloquantes : elles sont sous le seuil et marquées critiques, ce qui interdit un
passage en production tant qu'elles ne sont pas remédiées. Le tableau du bas montre que
chaque note est adossée à une preuve — c'est ce qui rend le diagnostic auditable.
\end{oral}

% =====================================================================
\section{Panneau de saisie du diagnostic}

Le bouton « Poursuivre le diagnostic » ouvre un panneau latéral listant les dimensions
de la campagne. Chaque ligne permet de régler la note, la confiance et la preuve. Une
dimension jamais instruite est explicitement signalée comme telle, et non affichée à
zéro.

L'enregistrement recalcule le readiness immédiatement : le score, les moyennes d'axe, le
taux de complétion et la liste des dimensions bloquantes sont rafraîchis sans
rechargement de page.

\capture{05-diagnostic-panneau.png}{Panneau de saisie des dimensions}{0.88}

\begin{oral}
Je modifie la note de la dimension « Gouvernance data » et je l'accompagne d'une preuve.
À l'enregistrement, le score global, les deux moyennes d'axe et la liste des dimensions
bloquantes sont recalculés côté serveur, puis réaffichés. Rien n'est stocké côté
navigateur : ce que vous voyez sort de la base.
\end{oral}

% =====================================================================
\section{Liste des opportunités}

L'écran recense les besoins métier. Un bandeau rappelle à l'utilisateur à périmètre
restreint qu'il ne voit que ses départements. Un second bandeau, en ambre, signale les
doublons probables détectés par le croisement interservices, en nommant les codes
concernés et les termes qu'ils partagent.

Le tableau se filtre par département et par recherche libre sur le titre ou le code.
Chaque ligne porte son type de solution recommandé, son score et ses actions. Le bouton
« Nouvelle opportunité » n'apparaît que si l'utilisateur détient
\code{opportunity.create}.

\capture{06-opportunites.png}{Liste des opportunités et détection de doublons}{0.95}

\begin{oral}
C'est ici que remontent les besoins. Le bandeau ambre est un des apports du croisement
interservices : deux départements ont exprimé un besoin très proche sans le savoir, et
la plateforme les rapproche en s'appuyant sur les termes partagés. Sans cet écran, les
deux projets seraient financés séparément.
\end{oral}

% =====================================================================
\section{Formulaire d'opportunité}

Le formulaire couvre trois blocs. Le rattachement d'abord — département et site — car le
département détermine le volet spécifique du questionnaire. Le contexte ensuite :
situation actuelle et points de douleur. Le questionnaire de qualification enfin, dont
les réponses alimentent directement les quinze faits du moteur.

À l'enregistrement, la qualification est déclenchée automatiquement. Il n'existe pas
d'opportunité non qualifiée dans la base.

\capture{07-opportunite-formulaire.png}{Création d'une opportunité}{0.92}

\begin{oral}
Le questionnaire n'est pas décoratif : chaque réponse est un fait d'entrée du moteur de
qualification. Je déclare ici des données disponibles mais non structurées, du texte
libre à traiter, et pas de génération de contenu. Je valide. Le moteur s'exécute
immédiatement et propose un type de solution. Notez qu'il n'y a pas d'étape « lancer la
qualification » : la création et la qualification sont une seule opération.
\end{oral}

% =====================================================================
\section{Architecture et données}

Cet écran tient le référentiel AS-IS en trois sections : les processus, les applications
du système d'information et les sources de données avec leur score de qualité et
l'indication d'une vérité terrain disponible.

C'est la matière première du moteur de qualification. Une source de données mal notée ou
dépourvue de vérité terrain fait échouer les règles bloquantes des types qui exigent un
apprentissage, et oriente mécaniquement la recommandation vers l'automatisation
déterministe ou la refonte de processus.

\capture{08-architecture.png}{Référentiel des processus, applications et données}{0.95}

\begin{oral}
Cet écran explique pourquoi la plateforme refuse parfois de recommander de l'IA. Si la
source de données d'un processus est notée 2 sur 5 et n'a pas de vérité terrain, la
règle bloquante du type « apprentissage supervisé » n'est pas satisfaite, et le moteur
bascule sur de l'automatisation ou sur une refonte de processus. La qualification
technique arrive donc avant la décision, pas après.
\end{oral}

% =====================================================================
\section{Portefeuille}

Le portefeuille présente toutes les opportunités classées par score, avec le bénéfice
attendu cumulé affiché en sous-titre. Cinq filtres sont disponibles : département, type
de solution, phase, score minimal et statut. Le tri est ajustable.

L'action « Recalculer la priorisation » n'est proposée qu'aux détenteurs de
\code{portfolio.prioritize}.

\capture{09-portefeuille.png}{Table du portefeuille priorisé}{0.95}

\begin{oral}
Le portefeuille est l'écran d'arbitrage. Chaque score est calculé sur six critères
pondérés, moins une pénalité de risque proportionnelle au produit probabilité par
impact. Le rang et la vague en découlent : la vague~1 est ce qu'on lance, la vague~2 ce
qu'on prépare. Rien n'est saisi à la main dans cette colonne.
\end{oral}

% =====================================================================
\section{Panneau de repriorisation}

Le panneau expose les six pondérations sous forme de curseurs. Toute grille saisie est
renormalisée à~1 avant application, de sorte qu'un utilisateur qui pousse un seul critère
obtient une pondération valide sans avoir à ajuster les cinq autres.

Le recalcul réévalue tout le portefeuille, réattribue les rangs et les vagues, et
journalise l'opération dans le journal d'audit avec l'ancienne et la nouvelle
pondération.

\capture{10-repriorisation.png}{Ajustement des pondérations de priorisation}{0.88}

\begin{oral}
C'est l'écran qui montre que la priorisation est un outil de pilotage, pas une formule
figée. Si la Direction décide que le contexte impose de privilégier l'impact financier,
elle pousse ce curseur et relance le calcul. Le classement change réellement : sur le
jeu de démonstration, douze opportunités changent de rang. Et le changement de
pondération est lui-même tracé dans le journal d'audit, avec l'ancienne et la nouvelle
grille.
\end{oral}

% =====================================================================
\section{Liste des fiches projet}

Un écran de sélection listant les fiches disponibles, chacune renvoyant au dossier de
cadrage correspondant. Les fiches ne sont pas des documents importés : elles sont
assemblées à la demande à partir des données vivantes du portefeuille.

\capture{11-fiches.png}{Sélection d'une fiche projet}{0.9}

\begin{oral}
Ces fiches ne sont pas stockées comme des documents. Elles sont reconstruites à chaque
ouverture à partir de l'opportunité, de sa recommandation, de son score, de ses
indicateurs, de ses risques et de son business case. Une fiche ne peut donc pas être
périmée par rapport aux données.
\end{oral}

% =====================================================================
\section{Fiche projet}

C'est l'écran le plus dense de l'application, et le point d'aboutissement de la
démarche. Un bandeau d'identité — département, sponsor, responsable métier, responsable
technique, phase et statut — précède deux colonnes.

La colonne principale déroule le contexte AS-IS et ses points de douleur, la solution
recommandée avec son patron de solution, sa justification rédigée, les règles satisfaites
affichées en pastilles, la première alternative écartée et les garde-fous à mettre en
place ; puis les indicateurs et le bénéfice attendu, et enfin les risques et la
conformité.

La colonne latérale, sur fond sombre, affiche le score de priorité en grand, le rang sur
la taille du portefeuille, la vague, puis la contribution détaillée de chacun des six
critères et la pénalité de risque. Viennent ensuite le business case, réservé aux
porteurs de \code{finance.read}, la décision de comité et un bloc d'explicabilité.

\capture{12-fiche-projet.png}{Fiche projet complète}{0.95}

\begin{oral}
Cette fiche est ce qu'un comité d'investissement reçoit. Trois choses la distinguent
d'un document classique. D'abord, la recommandation est justifiée : les pastilles vertes
sont les règles qui ont été satisfaites, et la pastille grise est l'alternative écartée
avec le motif. Ensuite, le score de la colonne de droite est décomposé critère par
critère, y compris la pénalité de risque. Enfin, les garde-fous sont générés
automatiquement — sur une opportunité en périmètre GxP, la plateforme ajoute la
validation informatique, la traçabilité ALCOA+ et la revue humaine obligatoire.
\end{oral}

\begin{ecart}
Le bouton porte le libellé « Export Word / PDF ». L'API produit en réalité un rendu JSON
et une page HTML mise en page pour l'impression : l'utilisateur obtient un PDF par la
fonction d'impression de son navigateur, mais aucun fichier \code{.docx} n'est généré.
Le libellé est en avance sur la fonction.
\end{ecart}

% =====================================================================
\section{Valeur et ROI}

L'écran consolide les business cases sur un horizon de trois ans. Il présente les gains
annuels par département, le classement des projets par ROI pour les huit premiers, le
suivi des indicateurs avant et après, et un tableau détaillé par projet donnant le coût
total de possession, les gains annuels, le ROI et le délai de retour.

\capture{13-valeur.png}{Pilotage de la valeur et du retour sur investissement}{0.95}

\begin{oral}
C'est l'écran qui ferme la boucle. Le business case ne sert plus seulement à obtenir un
budget : la colonne des indicateurs compare la valeur promise à la valeur effectivement
mesurée. Le délai de retour se calcule sur le seul investissement initial rapporté au
gain annuel net d'exploitation — y faire figurer trois années de coûts récurrents qui ne
sont pas encore engagés allongerait artificiellement le délai.
\end{oral}

% =====================================================================
\section{Administration}

L'écran d'administration est organisé en trois onglets, dont l'affichage dépend des
droits : \emph{Organisation}, \emph{Utilisateurs} et \emph{Rôles \& permissions}.

L'onglet Organisation regroupe les informations de l'entreprise, ses sites et ses
départements. L'onglet Utilisateurs permet de créer un compte, de l'activer ou de le
désactiver, de lui attribuer des rôles et de le rattacher à des départements et à des
sites. L'onglet Rôles affiche la grille rôle~$\times$~permission ; les rôles fournis par
la plateforme y sont en lecture seule, un rôle personnalisé devant être créé pour
ajuster les droits.

\capture{14-administration.png}{Console d'administration, onglet Rôles et permissions}{0.95}

\begin{oral}
Voici le module M1. Les onze rôles que vous voyez ne sont pas des libellés décoratifs :
chacun est un jeu de permissions parmi les quarante et une du catalogue, associé à un
niveau de périmètre. Les rôles système ne sont pas modifiables, pour qu'un
administrateur ne puisse pas se retirer involontairement des droits dont la plateforme
dépend. Il peut en revanche créer ses propres rôles. Une seule permission lui est
interdite : celle qui franchit la frontière entre organisations.
\end{oral}

% =====================================================================
\section{Journal d'audit}

Le dernier écran présente les événements en insertion seule. Chaque ligne porte
l'horodatage, l'auteur et son adresse électronique, l'action, l'entité concernée et son
libellé. Le détail expose l'ancienne et la nouvelle valeur.

Le sous-titre de l'écran énonce la propriété qui fait sa valeur : aucune entrée ne peut
être modifiée ni supprimée. L'API n'expose ni mise à jour ni suppression sur cette
table.

\capture{15-audit.png}{Journal d'audit en insertion seule}{0.95}

\begin{oral}
C'est l'écran qui répond à la question « pourquoi cette décision ? » six mois plus tard.
Chaque modification significative laisse une trace avec son auteur, sa date, l'ancienne
et la nouvelle valeur. Le point important est qu'il n'existe aucun moyen, ni par
l'interface ni par l'API, de modifier ou de supprimer une entrée : le journal est en
insertion seule. C'est une exigence directe du contexte réglementaire pharmaceutique.
\end{oral}
